\section{Phương pháp đề xuất}
\subsection{Tổ chức dữ liệu và ngân sách nhãn}
Chia dữ liệu ban đầu thành các tập tách biệt: $D_{\mathrm{train}}$ để huấn luyện, $D_{\mathrm{cal}}$ để calibration/fit ATC và $D_{\mathrm{ref}}$ để lấy mốc drift, chất lượng. Mỗi version có estimator và mốc chất lượng riêng; sau khi thay champion phải fit lại estimator từ calibration data hợp lệ và ghi nhận mốc mới.

Batch production $W_t$ được lưu prediction \emph{trước} khi tiết lộ nhãn. Trước khi đọc nhãn, dùng seed để phân vai các mẫu thành pool kiểm chứng, huấn luyện và đánh giá candidate. Các pool là mẫu ngẫu nhiên của cùng batch và không trùng ID. Chỉ pool kiểm chứng được dùng kết luận chất lượng champion; pool huấn luyện dùng tạo dataset; pool đánh giá dành riêng cho candidate. Nhãn không được tái sử dụng sang pool khác trong thí nghiệm chính.

Đặt $B$ là giới hạn số nhãn production được tiết lộ trên một episode. Ghi sổ ba khoản:
\[
B_{\mathrm{used}}=B_{\mathrm{verify}}+B_{\mathrm{train}}+B_{\mathrm{gate}}\le B.
\]
Tất cả nhãn trả về, kể cả khi không retrain hoặc candidate bị loại, đều được tính; retry cùng ID không tính lại. Quota theo mục đích và lịch cấp nhãn huấn luyện được cố định trên development. Trong so sánh policy kiểm chứng, các phương pháp có cùng lịch và pool nhãn huấn luyện, tránh lợi ích chỉ do được cấp thêm training labels. Quota kiểm chứng gồm phần dành cho audit định kỳ và phần cho yêu cầu bổ sung; không tiêu hết phần audit vào các cảnh báo sớm.

\subsection{Ước lượng và điều kiện xác nhận suy giảm}
Evidently theo dõi $X$ trên toàn batch hợp lệ. CBPE fit bằng reference/calibration có nhãn và ước lượng chất lượng từ xác suất lớp dương trên production~\cite{cbpe}. ATC được fit độc lập trên cùng split làm baseline~\cite{atc}. Cấu hình đầu tiên dùng \textbf{tỷ lệ lỗi $e=1-\mathrm{accuracy}$} làm metric quyết định chung vì cả hai estimator đều hỗ trợ; F1, recall và Brier là metric báo cáo bổ sung. Nếu lớp hiếm quá ít, không đưa kết luận riêng về recall của lớp đó.

Đặt $e_{\mathrm{ref}}$ là tỷ lệ lỗi champion trên $D_{\mathrm{ref}}$, $\Delta_t=e_t-e_{\mathrm{ref}}$ là mức suy giảm. Estimator cung cấp $\widehat{\Delta}_t$ để gợi ý lấy nhãn, không phải ground truth. Với khoảng kiểm chứng $I_t=[L_t,U_t]$ từ mẫu nhãn thật:
\begin{itemize}
  \item $L_t>\delta_{\mathrm{harm}}$: xác nhận suy giảm vượt biên thực tiễn.
  \item $U_t\le\delta_{\mathrm{keep}}$: có bằng chứng giữ mô hình trong biên chấp nhận.
  \item Các trường hợp khác: chưa xác định; với $0\le\delta_{\mathrm{keep}}\le\delta_{\mathrm{harm}}$.
\end{itemize}
Khoảng kiểm chứng phải tính bất định của cả mẫu production và mốc reference. Bản đầu cho phép tối đa hai lần đọc kết quả trong một cửa sổ, với cỡ mẫu dự kiến xác định trước. Lấy mẫu bằng tiền tố của một hoán vị ngẫu nhiên độc lập với nhãn; dùng cận cho sai số nhị phân và phân bổ mức lỗi thống kê giữa các lần đọc và mốc reference. Không dùng CI cố định rồi xem lặp vô hạn. Mức bảo đảm phát biểu theo cửa sổ; false-alarm toàn stream được đo riêng. Calibration và repeated testing là giới hạn cần báo cáo.

\subsection{Policy kiểm chứng trong giới hạn ngân sách}
Policy có tham số cố định: lịch audit $k$, cỡ mẫu ban đầu $m_1$, cỡ mẫu bổ sung $m_2$, ngưỡng drift, ngưỡng estimator, mức bất đồng và cooldown. Thực hiện theo thứ tự:
\begin{enumerate}
  \item Kiểm tra contract, tính đầy đủ của batch và reference. Nếu lỗi, trả HOLD và ghi lý do; chưa dùng batch để retrain.
  \item Tính drift và ước lượng chất lượng. Bắt đầu kiểm chứng khi đến kỳ audit, có drift hoặc estimator nghi suy giảm. Audit định kỳ vẫn chạy khi cả hai tín hiệu im lặng.
  \item Nếu còn quota, yêu cầu mẫu ngẫu nhiên $m_1$ ở pool kiểm chứng. Khi $I_t$ đã phân loại được, quyết định theo bảng~\ref{tab:policy}.
  \item Nếu chưa xác định hoặc estimator bất đồng rõ với nhãn thật, dùng tối đa một lượt bổ sung $m_2$ từ cùng hoán vị, rồi cập nhật khoảng kiểm chứng theo quy tắc hai lượt. Estimator không có quyền phủ quyết kết quả nhãn thật.
  \item Khi hết quota, thiếu mẫu hoặc vẫn chưa xác định, trả HOLD, ghi hẹn audit tiếp theo. Không diễn giải HOLD thành mô hình khỏe.
\end{enumerate}

\begin{table}[ht]
\centering
\caption{Quyết định sau kiểm chứng.}\label{tab:policy}
\begin{tabularx}{\linewidth}{@{}L{3.5cm} Y@{}}
\toprule\textbf{Hành động} & \textbf{Điều kiện và kết quả} \\
\midrule
KEEP & Khoảng kiểm chứng nằm trong biên chấp nhận; lưu bằng chứng, tiếp tục giám sát. \\
REQUEST\_LABELS & Cần kiểm chứng và còn lượt/quota; chọn ID trước khi đọc nhãn, lưu chi phí. \\
HOLD & Chưa đủ bằng chứng, hết ngân sách, dữ liệu lỗi hoặc chưa đủ dữ liệu huấn luyện; có lý do và thời điểm xét lại. \\
RETRAIN & Nhãn thật xác nhận suy giảm, pool training đủ mẫu/lớp, còn quota đánh giá và không có job cùng champion đang chạy. Tạo candidate theo recipe cố định. \\
\bottomrule
\end{tabularx}
\end{table}
Phần cần kiểm chứng thực nghiệm là lợi ích của phân bổ audit theo bằng chứng so với lịch cố định. Không tự động học hay thay đổi tham số policy trên test. Không trộn F1 đo từ mẫu có nhãn với F1 ước lượng trên phần không nhãn bằng trung bình tùy ý.

\subsection{Recipe retraining và đánh giá candidate}
Recipe giữ nguyên model family, preprocessing, hyperparameters và phương thức huấn luyện từ đầu. Dataset gồm nhãn production thuộc pool training trong cửa sổ gần nhất, trộn tỷ lệ cố định từ $D_{\mathrm{train}}$. Dùng lấy mẫu ngẫu nhiên/phân tầng với seed; cố định cửa sổ, tỷ lệ trộn và giới hạn số mẫu qua pilot. Dành calibration split không trùng training cho candidate. Khi thiếu dữ liệu/lớp, HOLD thay vì dùng nhãn giả.

Luồng thực thi là \textbf{TrainingJob $\rightarrow$ TrainingOutput $\rightarrow$ Build/Register candidate $\rightarrow$ EvaluationGate}. Gate dùng mẫu production riêng có nhãn được tính vào $B_{\mathrm{gate}}$; so sánh lỗi candidate và champion theo cặp trên cùng mẫu. Candidate chỉ pass khi lợi ích vượt biên đã định và không suy giảm quá biên trên historical holdout. Một gate dùng cỡ mẫu định trước; không mở rộng mẫu tùy kết quả trong MVP. Thiếu bằng chứng trả HOLD; fail không đổi champion. Mọi candidate đều được tính chi phí.

Trong replay, candidate pass được kích hoạt cho batch sau khi job/gate hoàn tất; champion cũ phục vụ trong thời gian chờ. Trên AWS/K3s, luồng Argo thực hiện cập nhật có kiểm soát và kiểm tra phiên bản phục vụ thực tế theo Mục~\ref{sec:orchestration}. Không đồng nhất đăng ký version với cải thiện chất lượng hay quyền tự động triển khai: chỉ cập nhật khi gate và các điều kiện vận hành đều đạt.
